\section{Conclusões e trabalhos futuros}
\label{sec:conclusoes}
A comparação de AE e IF em 523.937.952 transações da Ethereum mostrou que mecanismos distintos de anomalia geram conjuntos complementares e respostas temporais diferentes, mesmo com atributos e limiares de calibração comuns ao protocolo. Com hiperparâmetros preselecionados e novos ajustes em 2025, o IF apresentou maior AP em todos os recortes. No teste de julho a novembro, recuperou 757.192 transações associadas a bots MEV em Q99,5, frente às 544.738 do AE, com menor volume sinalizado. A interseção concentrou 36,70\% de associadas entre 956.106 sinalizações, enquanto a união ampliou a recuperação para 951.073 associadas.

Em dezembro, o AE passou a contribuir com a maior parcela exclusiva da união associada. No estresse, recuperou 75.354 associadas em Q99,5, contra 13.313 do IF, e concentrou 7,12\% no topo de 1\%, frente a 1,71\%. A AP do IF permaneceu maior nessa janela, distinguindo a ordenação global da concentração obtida sob um orçamento específico. Os achados sustentam a comparação por regime e por objetivo de investigação: interseção, união e ranking individual oferecem compromissos mensuráveis entre concentração, recuperação e esforço de análise.

Como continuidade, propomos adaptar os algoritmos de Torres et al.\ \cite{Torres2021} ao contexto da Ethereum de 2025 para identificar candidatos a padrões de inserção, deslocamento e supressão. A geração estrutural será independente dos escores e fornecerá uma referência comum para comparar a cobertura do AE, do IF e de suas sinalizações compartilhadas e exclusivas.

Os candidatos serão submetidos a avaliação semântica, articulando ordem transacional, chamadas de contratos, recibos, logs, rastros de execução e fluxos de ativos. A investigação examinará os papéis dos participantes e os efeitos econômicos, com simulações de cenários alternativos de execução quando pertinentes. Essa extensão permitirá validar padrões concretos de \textit{front-running} e medir como cada política de priorização recupera eventos confirmados. A integração entre escores não supervisionados, detecção estrutural e evidência semântica amplia o uso dos resultados como instrumento de investigação da atividade transacional.
